\section{Permuted split alternating constituents with a complete partner}
\label{sec:permuted-split-partner}

\paragraph{Checked status and provenance.}
HD-TW06--07, question Q2dw, is a root extension of the accepted
HD-TW04--05 weighted argument and HD-FT03 split-character estimate.
The author and independent reviewer have read the complete frozen proof and source interfaces; both whole-proof checks pass without repair. No numerical
certificate is required or claimed. The source and complete-partner
hypotheses below are actual representation hypotheses.

\subsection*{HD-TW06. The cyclic local and weighted arguments only need compatible products}

Let \(H\subset S_n\), and let \(W\) be a complex representation of
\(H\), of dimension \(a\). Suppose
\[
 n\ge24,\quad a\ge n^3,\qquad
 \frac{|\chi_W(y)|}{a}\le
       \left(\frac{\max\{f(y),1\}}n\right)^3
       \quad(y\in H),
 \tag{TWS.1}
\]
where \(f(y)\) counts fixed natural letters. Define
\[
 A=\frac{a(n-4)}{n^2},\qquad B=\frac an,
 \qquad L=a-(n-1)A=\frac{a(5n-4)}{n^2},
 \qquad \lambda=1-\frac{nA+1}{a}=\frac4n-\frac1a.
 \tag{TWS.2}
\]
Then the purely local conclusion of HD-TW02 holds for every
nonidentity \(x\in H\):
\[
                 \nu(W(x))\ge A\operatorname{ind}(x)+B.
                 \tag{TWS.3}
\]
This statement needs no extension of \(W\) to \(S_n\).

\begin{proof}
Write \(I=\operatorname{ind}(x)\), \(q=\operatorname{ord}(x)\),
and \(f_j=f(x^j)\). Every power lies in \(H\), so cyclic Fourier
inversion and the triangle inequality give, for every target eigenvalue,
\[
 \frac{\nu(W(x))}{a}\ge
  1-\frac1q\sum_{j=0}^{q-1}
        \left(\frac{\max\{f_j,1\}}n\right)^3,
 \qquad \frac1q\sum_jf_j=n-I.
 \tag{TWS.4}
\]
The second identity counts cyclic orbits on the letters. If
\(I\ge n/3\), then
\((\max\{f,1\}/n)^3\le f/n+n^{-3}\) for \(0\le f\le n\).
Subtracting \(AI/a\) from (TWS.4) leaves at least
\(4I/n^2-n^{-3}\ge4/(3n)-n^{-3}>1/n\).
If \(I<n/3\), the moved support of \(x\) has at most \(2I\)
letters, since each nontrivial cycle of length \(d\) contributes
\(d-1\ge d/2\) to its index. Hence all powers have
\(u_j=f_j/n>1/3\). The inequality
\(1-u_j^3\ge(13/9)(1-u_j)\) then shows that the same difference
is at least \(4I/(9n)+4I/n^2>1/n\) when \(I\ge3\).
For the remaining indices \(I=1,2\), one has
\(u_j\ge1-4/n\ge5/6\) and
\(1+u_j+u_j^2\ge91/36>5/2\). The difference is therefore at
least \(3I/(2n)+4I/n^2>1/n\). These cases prove (TWS.3).
All operators and all powers used belonged to \(H\), which explains
the subgroup scope of this form of the accepted local lemma.
\end{proof}

Now assume an actual connected finite-group sphere cover has a genuine
generating product-one tuple \(T_1,\ldots,T_r\), with
\[
 V=W^{\otimes t}\otimes U,\quad t\ge2,\quad b=\dim U\ge1,
 \quad e=a^t b,\quad r\ge12,
 \qquad\sum_i\operatorname{rank}(1-T_i)=2e+6.
 \tag{TWS.5}
\]
The whole first block is stable and its factors are permuted
transitively. Its actual action on the factor algebras supplies
compatible cycle products \(W(x)\), up to scalars, with \(x\in H\).
The factor permutation on \(t\) letters and the associated block
permutation on \(nt\) natural letters are actual transitive quotient
actions of this cover. Suppose the entire partner satisfies
\[
        \nu(S)/b\ge1/n\quad
        \hbox{for every nonscalar operator in its projective image}.
        \tag{TWS.6}
\]
For \(b=1\) this is vacuous. Under (TWS.1), no such source exists.

\begin{proof}
We give the weighted extension explicitly to retain the subgroup
interface. At a factor-fixing entry with at least one nonidentity
within-factor permutation, spectral convolution and averaging over
the active factors give normalized moving rank at least
\(AI_i/(ta)+B/a\), where \(I_i\) is its natural block index.
If every within-factor permutation is identity, the first block is
scalar. A nonscalar entire partner gives \(1/n=B/a\) by (TWS.6),
whereas a scalar partner gives full rank by genuineness of the
original entry. Thus the same additive bound holds for every
factor-fixing entry, with no projected mark-count assumption.

The accepted weighted-word cycle bound of D-ATT08 gives every
factor-moving entry normalized moving rank at least
\((a-1)/(2a)\). Since \(2B\le a-1\), every genuine entry
has moving rank at least \(Be/a\). Four factor-moving entries
and eight other entries would therefore give
\[
 D\ge2e+(8B-2)e/a>2e+6,
 \tag{TWS.7}
\]
because \(e/a\ge a\) and \(a(8B-2)>6\).
Consequently at most three entries move factors and at least nine
fix them.

For any entry let \(\pi_i\) be its factor permutation, with cycle
lengths \(\ell_j\) and compatible cycle products \(x_j\in H\).
The same accepted exact short-word estimate gives each cycle's
normalized support at least
\[
 1-\frac1{\ell_j}
   +\frac{\nu(W(x_j))}{\ell_j a}
   -\frac{R_{\ell_j}(a)}{\ell_j a^{\ell_j}},
 \quad R_1=0,\quad
 R_{\ell}(a)/a^\ell\le1/a\quad(\ell\ge2).
 \tag{TWS.8}
\]
This estimate concerns an actual tensor cycle and every scalar target;
it does not require a symmetric-group extension. Weighting its bounds
by \(\ell_j/t\), and using at most \(\operatorname{ind}(\pi_i)\)
nontrivial cycles, gives
\[
 \frac{\operatorname{rank}(1-T_i)}e
  \ge\frac{AI_i}{ta}
       +\frac{\lambda\operatorname{ind}(\pi_i)}t,
 \qquad I_i=n\operatorname{ind}(\pi_i)+
                         \sum_j\operatorname{ind}(x_j).
 \tag{TWS.9}
\]
Here \(\nu(W(x))\ge A\operatorname{ind}(x)\) is also valid
at identity products. The additional \(B/a\) remains available
at the factor-fixing entries only. The actual two quotient covers
give \(\sum I_i\ge2nt-2\) and
\(\sum\operatorname{ind}(\pi_i)\ge2t-2\).
Since \(A,\lambda\ge0\), their sum yields
\[
 D-2e\ge\frac ea
       \left(r_{\rm fix}B-2-\frac2t(L-1)\right)
 \ge\frac ea(9B-L-1)>6.
 \tag{TWS.10}
\]
For completeness, all numerical conditions used follow directly from
\(n\ge24\), \(a\ge n^3\). One has \(A\ge0\), \(B>0\),
\(2B\le a-1\), \(L\ge1\), \(\lambda>0\), and
\[
 8B-2=8a/n-2,\qquad
 9B-L-1=\frac{4a(n+1)}{n^2}-1.
 \tag{TWS.11}
\]
Both expressions multiplied by \(a\) exceed six; for the first
use \(a/n\ge n^2\), and for the second use
\(a(n+1)/n^2\ge n(n+1)\).
This contradicts (TWS.5) and proves the assertion.
\end{proof}

\subsection*{HD-TW07. The nonextendible ordinary split case}

\begin{theorem}[HD-TW07]
In (TWS.5), suppose the actual stable core block is an irreducible
module for \(A_n^t\), whose simple factors are permuted transitively
by the coefficient image. Suppose each factor module is a split
ordinary alternating constituent, rather than an extendible
ordinary module. Retain the actual natural quotient actions and
entire-partner hypothesis (TWS.6). If
\[
                              n\ge65536,
                              \tag{TWS.12}
\]
then the source is impossible, for every \(t\ge2\).
No extension of either split constituent to \(S_n\) is asserted.
\end{theorem}

\begin{proof}
For \(n\ge7\), every automorphism of \(A_n\) comes from
conjugation by a natural permutation. An odd permutation exchanges
the two split constituents of a self-conjugate label, while inner
automorphisms preserve each. These are precisely the accepted
automorphism and split-character inputs in the compatible-normalizer
paragraph of HD-FT03.

Transitivity on the simple factors therefore makes their modules
belong to the same self-conjugate label \(\lambda\), with a choice
of one of its two constituents at each factor. At any factor with
the other constituent, change the identification of its natural
\(A_n\) by conjugation with an odd permutation. After these choices
all factor modules identify with one fixed constituent \(W\).
This only conjugates the natural block action by a fixed permutation
inside each block, so indices, transitivity, and actual quotient
status are unchanged.

Every compatible local intertwiner after this change implements an
inner automorphism. Indeed an odd automorphism would interchange the
two inequivalent constituents. It cannot even preserve the projective
class of \(W\): an equality up to scalar would give a linear
character of the perfect group \(A_n\), hence an equality of ordinary
representations, which is impossible. An intertwiner for an inner
automorphism is a scalar times its representing matrix by Schur's
lemma. In particular every compatible product around a factor cycle
has the form \(W(x)\) up to scalar, with \(x\in A_n\).

The full accepted estimate HD-FT03, equation (FTS.1), states at
\(n\ge65536\) that \(a=\dim W\ge n^{120}\) and
\[
 \frac{|\chi_W(x)|}{a}\le
 \begin{cases}
 (f(x)/n)^3,&f(x)>0,\\
 n^{-6},&f(x)=0.
 \end{cases}
 \tag{TWS.13}
\]
Since \(n^{-6}\le n^{-3}\), this is (TWS.1) for
\(H=A_n\). The compatible products just established satisfy the
entire local interface of HD-TW06, and \(a\ge n^3\).
That theorem gives the contradiction. This proof uses the published
split-character and effective MSP inputs only through the completely
proved HD-FT03 estimate; their exact source and numerical-sketch
boundaries remain those recorded there.
\end{proof}

\paragraph{Scope and geometric meaning.}
This closes the permuted ordinary split branch under the entire-partner
floor, a branch not covered by the stated extendible-ordinary scope of
HD-TW05. A single stable split factor is separately excluded by
HD-FT03 for \(n>3500000\), with arbitrary partners and no projected
mark-count hypothesis. Neither result covers spin factors or selects
an individual quotient from a permuted orbit. No source, polarization,
cycle, or detector is inferred from these exclusions.

In the programme's terms, this removes a possible curve supplier for
condition A, cote16 scan050 as corrected in SIGN\_NEXT SN-C03/C04.
It does not supply condition I positivity, scan049 and cote44 Theorem4,
nor replace the Cayley--Hamilton projector of scans048--049 or the
previously stated SGA/EGA/FGA spreading interfaces. Markman, Severn
and032 are not premises.
